% GENERATED FILE. Edit canonical lessons, then run npm run generate:books.
% canonical-source-hash: fnv1a64:b9dc02df75141be8
% canonical-lessons: RU-C122-chin, RU-C122-rainbow, RU-C122-lightning, RU-C122-thunder, RU-C122-planet

\chapter{Chin, Rainbow, Lightning, Thunder, Planet}
\label{ch:ru-chin-rainbow-lightning-thunder-planet}
\hlchaptermodality{122}

\begin{chapteropening}
\textbf{By the end of this chapter:} \emph{I can say a chin, a rainbow, lightning, thunder and a planet.}

Complete the last lesson of chapter 122: I can say a chin, a rainbow, lightning, thunder and a planet.
\end{chapteropening}

\section[\texorpdfstring{\ru{подбородок} (podboródok)}{подбородок (podboródok)}]{\ru{подбородок} (podboródok) — a chin}

\label{lesson:RU-C122-chin}



\begin{quote}
Before the new one: say the Russian for a stomach, then the Russian for a throat.
\end{quote}

\subsection*{You'll want to know: \ru{подбородок}}
\textbf{\ru{подбородок}} (\emph{podboródok}) — \textquotedblleft{}a chin\textquotedblright{}.

\begin{cousinweb}
\textbf{\ru{подбородок}} is \emph{\ru{под}}, \textquotedblleft{}under\textquotedblright{}, with \emph{\ru{борода}}, \textquotedblleft{}beard\textquotedblright{}, and \emph{-\ru{ок}}: the part under the beard. Vasmer relates \emph{\ru{борода}} to Lithuanian \emph{barzdà}, Latin \emph{barba} and Old High German \emph{bart}: English \emph{beard} is a cousin.
\end{cousinweb}

\begin{grammarlens}[title={What you've built}]
One more word for this chapter.
\end{grammarlens}

\subsection*{Your turn}
\begin{itemize}
\raggedright
  \item \emph{Say it:} \emph{podboródok}
  \item \emph{Say it:} \emph{podboródok}, once more
  \item \emph{Say it:} say \emph{podboródok}
  \item \emph{Recall:} say the Russian for a stomach, then the Russian for a throat, then say \emph{podboródok} again
\end{itemize}

\begin{tcolorbox}[breakable,colback=teal!4,colframe=teal!35!black,title={Before you move on}]
What does \ru{подбородок} mean? (\textquotedblleft{}A chin\textquotedblright{}.) Say it once more.
\end{tcolorbox}
\section[\texorpdfstring{\ru{радуга} (ráduga)}{радуга (ráduga)}]{\ru{радуга} (ráduga) — a rainbow}

\label{lesson:RU-C122-rainbow}



\begin{quote}
Before the new one: say the Russian for a throat, then the Russian for a chin.
\end{quote}

\subsection*{You'll want to know: \ru{радуга}}
\textbf{\ru{радуга}} (\emph{ráduga}) — \textquotedblleft{}a rainbow\textquotedblright{}.

\begin{cousinweb}
\textbf{\ru{радуга}} — Vasmer finds most likely a derivation \emph{*radǫga} from \emph{\ru{рад}}, \textquotedblleft{}glad\textquotedblright{}, since dialects also call the rainbow \emph{\ru{весёлка}}, \textquotedblleft{}the merry one\textquotedblright{}. The form \emph{\ru{райдуга}}, with \emph{\ru{рай}}, \textquotedblleft{}paradise\textquotedblright{}, he takes as folk etymology.
\end{cousinweb}

\begin{grammarlens}[title={What you've built}]
One more word for this chapter.
\end{grammarlens}

\subsection*{Your turn}
\begin{itemize}
\raggedright
  \item \emph{Say it:} \emph{ráduga}
  \item \emph{Say it:} \emph{ráduga}, once more
  \item \emph{Say it:} say \emph{ráduga}
  \item \emph{Recall:} say the Russian for a throat, then the Russian for a chin, then say \emph{ráduga} again
\end{itemize}

\begin{tcolorbox}[breakable,colback=teal!4,colframe=teal!35!black,title={Before you move on}]
What does \ru{радуга} mean? (\textquotedblleft{}A rainbow\textquotedblright{}.) Say it once more.
\end{tcolorbox}
\section[\texorpdfstring{\ru{молния} (mólniya)}{молния (mólniya)}]{\ru{молния} (mólniya) — lightning}

\label{lesson:RU-C122-lightning}



\begin{quote}
Before the new one: say the Russian for a chin, then the Russian for a rainbow.
\end{quote}

\subsection*{You'll want to know: \ru{молния}}
\textbf{\ru{молния}} (\emph{mólniya}) — \textquotedblleft{}lightning\textquotedblright{}.

\begin{cousinweb}
\textbf{\ru{молния}} $\leftarrow$ Proto-Slavic \emph{*mŭlnĭji} — inherited. Vasmer relates it to Old Prussian \emph{mealde}, \textquotedblleft{}lightning\textquotedblright{}, Welsh \emph{mellt}, \textquotedblleft{}lightning\textquotedblright{}, and Old Norse \emph{Mjǫllnir}, the name of Thor's hammer.
\end{cousinweb}

\begin{grammarlens}[title={What you've built}]
One more word for this chapter.
\end{grammarlens}

\subsection*{Your turn}
\begin{itemize}
\raggedright
  \item \emph{Say it:} \emph{mólniya}
  \item \emph{Say it:} \emph{mólniya}, once more
  \item \emph{Say it:} say \emph{mólniya}
  \item \emph{Recall:} say the Russian for a chin, then the Russian for a rainbow, then say \emph{mólniya} again
\end{itemize}

\begin{tcolorbox}[breakable,colback=teal!4,colframe=teal!35!black,title={Before you move on}]
What does \ru{молния} mean? (\textquotedblleft{}Lightning\textquotedblright{}.) Say it once more.
\end{tcolorbox}
\section[\texorpdfstring{\ru{гром} (grom)}{гром (grom)}]{\ru{гром} (grom) — thunder}

\label{lesson:RU-C122-thunder}



\begin{quote}
Before the new one: say the Russian for a rainbow, then the Russian for lightning.
\end{quote}

\subsection*{You'll want to know: \ru{гром}}
\textbf{\ru{гром}} (\emph{grom}) — \textquotedblleft{}thunder\textquotedblright{}.

\begin{cousinweb}
\textbf{\ru{гром}} $\leftarrow$ Proto-Slavic \emph{*gromŭ} — inherited, with \emph{\ru{греметь}}, \textquotedblleft{}to rumble\textquotedblright{} (Vasmer). Vasmer relates it to Old Prussian \emph{grumins}, \textquotedblleft{}thunder\textquotedblright{}, Lithuanian \emph{gruménti}, \textquotedblleft{}to rumble\textquotedblright{}, and Greek \emph{chrómos}, \textquotedblleft{}a crash\textquotedblright{}.
\end{cousinweb}

\begin{grammarlens}[title={What you've built}]
One more word for this chapter.
\end{grammarlens}

\subsection*{Your turn}
\begin{itemize}
\raggedright
  \item \emph{Say it:} \emph{grom}
  \item \emph{Say it:} \emph{grom}, once more
  \item \emph{Say it:} say \emph{grom}
  \item \emph{Recall:} say the Russian for a rainbow, then the Russian for lightning, then say \emph{grom} again
\end{itemize}

\begin{tcolorbox}[breakable,colback=teal!4,colframe=teal!35!black,title={Before you move on}]
What does \ru{гром} mean? (\textquotedblleft{}Thunder\textquotedblright{}.) Say it once more.
\end{tcolorbox}
\section[\texorpdfstring{\ru{планета} (planéta)}{планета (planéta)}]{\ru{планета} (planéta) — a planet}

\label{lesson:RU-C122-planet}



\begin{quote}
Before the new one: say the Russian for lightning, then the Russian for thunder.
\end{quote}

\subsection*{You'll want to know: \ru{планета}}
\textbf{\ru{планета}} (\emph{planéta}) — \textquotedblleft{}a planet\textquotedblright{}.

\begin{cousinweb}
\textbf{\ru{планета}} came through Polish \emph{planeta} from Latin \emph{planēta}, Greek \emph{plan\'{\={e}}tēs}, \textquotedblleft{}a wanderer\textquotedblright{}: a planet is a wandering star (Vasmer).
\end{cousinweb}

\begin{grammarlens}[title={What you've built}]
One more word for this chapter.
\end{grammarlens}

\subsection*{Your turn}
\begin{itemize}
\raggedright
  \item \emph{Say it:} \emph{planéta}
  \item \emph{Say it:} \emph{planéta}, once more
  \item \emph{Say it:} say \emph{planéta}
  \item \emph{Recall:} say the Russian for lightning, then the Russian for thunder, then say \emph{planéta} again
\end{itemize}

\begin{tcolorbox}[breakable,colback=teal!4,colframe=teal!35!black,title={Before you move on}]
What does \ru{планета} mean? (\textquotedblleft{}A planet\textquotedblright{}.) Say it once more.
\end{tcolorbox}
